%% ma: L12-B07-0023
%% lop: 12
%% chuong: 2
%% bai: 7
%% dang: 12.7.6
%% loai: TN4
%% muc: NB
%% dap_an: "A"
%% nguon: "CĐ2.2.pdf giáo viên gửi 27/10/2026 (M07, Lê Minh Khôi – Nguyễn Anh Khôi, Chương II Chủ đề 2 Tọa độ của vectơ trong không gian), Câu 20"
%% kien_thuc: "Bài 7 (hình chiếu của điểm lên mặt phẳng tọa độ)"
%% trang_thai: chinh-thuc
\begin{ex}
Trong không gian $Oxyz$, cho điểm $U\left(\dfrac32;\dfrac1{\sqrt2};-\dfrac3{\sqrt2}\right)$. Gọi $U_1$ là điểm đối xứng của $U$ qua mặt phẳng $(Oxy)$. Tọa độ của điểm $U_1$ là
\choice{\True $U_1\left(\dfrac32;\dfrac1{\sqrt2};\dfrac3{\sqrt2}\right)$}{$U_1\left(-\dfrac32;-\dfrac1{\sqrt2};-\dfrac3{\sqrt2}\right)$}{$U_1\left(-\dfrac32;-\dfrac1{\sqrt2};\dfrac3{\sqrt2}\right)$}{$U_1\left(\dfrac32;-\dfrac1{\sqrt2};\dfrac3{\sqrt2}\right)$}
\loigiai{Đối xứng qua $(Oxy)$: đổi dấu cao độ: $U_1\left(\dfrac32;\dfrac1{\sqrt2};\dfrac3{\sqrt2}\right)$.}
\end{ex}
